\begin{textsample}{POS Dim 3 – rj – Score 14.00 – rj/rj\_inf\_cam\_exp\_9.txt}  \label{ex:f3_pos_014}
Direitos de Aprendizagem no Campo de Experiências Traços, \textbf{Sons}, Cores e Formas - CONVIVER e fruir com os \textbf{colegas} e professores manifestações artísticas e culturais da sua comunidade e de outras culturas - artes \textbf{plásticas}, música, dança, teatro, cinema, folguedos e festas populares. - \textbf{BRINCAR} com diferentes \textbf{sons}, ritmos, formas, cores, \textbf{texturas}, objetos, materiais, construindo cenários e indumentárias para \textbf{brincadeiras} de faz-de-conta, encenações ou para festas tradicionais. - EXPLORAR variadas possibilidades de usos e combinações de materiais, substâncias, objetos e recursos tecnológicos para criar desenhos, modelagens, músicas, danças, encenações teatrais e musicais. - PARTICIPAR de decisões e ações relativas à organização do ambiente ( tanto o cotidiano quanto o preparado para determinados eventos ), à definição de temas e à escolha de materiais a serem usados em atividades \textbf{lúdicas} e artísticas. - EXPRESSAR suas emoções, sentimentos, necessidades e ideias cantando, dançando, esculpindo, desenhando, encenando. - CONHECER -SE no contato criativo com manifestações artísticas e culturais locais e de outras comunidades Práticas que auxiliam na garantia dos Direitos de Aprendizagem no Campo de Experiências Traços, \textbf{Sons}, Cores e Formas A configuração dos campos de experiências assegura a interação de diferentes linguagens, como ocorre quando as \textbf{crianças} criam um desenho a partir da audição de algumas músicas ou criam uma música a partir de \textbf{pinturas} ou desenhos.
Também pode ser planejada a integração da música e das linguagens visuais com a dança, o teatro, e com a literatura, acompanhando as \textbf{crianças} na narrativa de histórias com objetos sonoros e instrumentos musicais, pintando um cenário para a dramatização de uma história por elas \textbf{inventada}, dentre outras possibilidades.
Esse campo de experiências chama a atenção das professoras para a importância de educar a sensibilidade da \textbf{criança} por uma ação que seja ao mesmo tempo política, estética e ética, de incentivá -la a ter um agir \textbf{lúdico} e um olhar poético sobre o mundo e as pessoas e coisas nele existentes, de ampliar a percepção pela \textbf{criança} de cores, \textbf{sons}, silêncios, \textbf{texturas}, tamanhos, sabores, cheiros, a partir de sua corporeidade.
A sonoridade e a visualidade tornam -se conquistas das \textbf{crianças} quando elas participam de ambientes onde o prazer estético abre possibilidades.
Diante disso importa verificar quais aspectos do ambiente da unidade de Educação \textbf{Infantil} atuam como recursos para elas sentirem, explorarem, representarem, imaginarem, criarem.
Tais aspectos incluem desde as condições de tempo, espaço e materiais que lhes são disponibilizados, como as relações de estímulo e confianças que estabelecem com pessoas presentes, pois os vínculos estabelecidos entre as \textbf{crianças}, a \textbf{professora} e a cultura em que estão imersos criam motivos e incentivos para elas explorarem o ambiente, reconhecerem seus aspectos significativos e os expressarem de diferentes formas.
Ao propiciar às \textbf{crianças} experiências com a surpresa, a alegria, o questionamento, a descoberta e o encantamento, o olhar sensível da \textbf{professora} acompanha as muitas formas pelas quais elas se interrogam sobre o mundo e sobre si, trilham universos simbólicos presentes em sua cultura e em outras, e imergem em situações diversas onde se \textbf{emocionam} com o belo.

% matched lemmas: brincadeira, brincar, colega, criança, emocionar, infantil, inventar, lúdico, pintura, plástico, professora, som, textura
\end{textsample}
